% page 424 of the facsimile (image p-423 of the scan)
% block 147.3 x 242.0 mm ; left margin 8.0 mm
\pgc{-12.700mm}{0.000mm}{
\l{\cel{3}416.}
\l{}
\l{\sou{paketboto}.\cel{2}(\sou{navig.})\cel{2}Navo extreme kapaca, movata per}
\l{\cel{3}aquo-vaporo, petrolo(e mem per atom-forco) qua trans-}
\l{\cel{3}portas, trans la oceani, pasajanti, vari, postalaji,}
\l{\cel{3}e c. DEFIRS.}
\l{}
\l{\sou{pakidermo}. (\sou{zool.)}\cel{2}Klaso de mamiferi ne-ruminera, di}
\l{\cel{3}qui la pelo esas tre dika (elefanto, rinoceroso, e c.)}
\l{\cel{3}- DEFIS.}
\l{}
\l{\sou{pakotilio}. I.Vari ne-grava quin la maristi di navo darfas}
\l{\cel{3}kunportar por komercor, sen pagar freto-preco. - II.}
\l{\cel{3}Sortaro de vari destinita a kambieso, a komerco, en la}
\l{\cel{3}landi fora. - DFIS.}
\l{}
\l{\sou{paktar}.\cel{2}(\sou{netrans., kun)}.\cel{2}Agar konvenciono solena. -DEFIS}
\l{}
\l{\sou{pala.}\cel{2}I. (\sou{aludante la vizajo)}. Senkolora, terne blanka. -}
\l{\cel{3}II. Qua esas poke lumoza, poke koloroza. - EFIS.}
\l{}
\l{\sou{palaco}. Habiteyo luxoza di rejo, di princo. Domp splen-}
\l{\cel{3}dida di richa privato. - DEFIRS.}
\l{}
\l{\sou{palado}. (\sou{kemio).}\cel{2}Metalo blanka, tre maleebla, poke gi-}
\l{\cel{3}sebla. -\cel{1}\sou{Simbolo kemiala}\cel{1}:\cel{2}Pd.}
\l{}
\l{\sou{paladio}. (en\cel{1}\sou{la epoki antiqua).}\cel{2}Statuo de Pallas, e}
\l{\cel{3}gardita da la Troyani kom salvanto di lia urbo. - II.}
\l{\cel{3}To quon populo egardas kom sekuriganta sua salveso.}
\l{\cel{3}- DEFIRS.}
\l{}
\l{\sou{paladino}.\cel{2}Heroo di la +chevalier-epoko. -\cel{1}\sou{Anke metaf}.}
\l{\cel{3}DEFIRS.}
\l{}
\l{\sou{palankino}. En India, Chinia (ante Mao Tze Tung), e c.,}
\l{\cel{3}porto-lito quan servisti portas de sua shultro.}
\l{\cel{3}(Anke en la Egiptia faraonala). - DEFIS.}
\l{}
\l{\sou{palato}. (\sou{anat.)}\cel{2}Parto supra di la boko-kavajo. - EFIS.}
\l{}
\l{\sou{palateno}. (en\cel{1}\sou{la epoki feudala)}. Qua havas ofico en la}
\l{\cel{3}palaco di suvereno. - DEFIRS.}
\l{}
\l{\sou{palatino}. Furo per qua la mulieri kovras sua shultri.}
\l{\cel{3}- DFIS.}
\l{}
\l{\sou{paleografio}. Cienco-fako di la persono qua su okupas}
\l{\cel{3}dechifrar la surskriburi, la manuskripti antiqua o tre}
\l{\cel{3}olda, la charti, la diplomi, e c. - DEFIRS.}
\l{}
\l{\sou{paleontologio.}\cel{1}Historio naturala pri la animali e pri}
\l{\cel{3}la vejetinti fosila. - DEFIRS.}
\l{}
\l{paleozoologio. Parto di paleontologio di qua la studiajo}
\l{\cel{3}esas la animali fosila. - DEFIS.}
}
